\subsection{Five at-most-one encodings}
Let $X=(x_1,\ldots,x_k)$ be an ordered sequence of distinct input variables.
An encoding $E(X,Y)$ may introduce auxiliary variables $Y$, but must satisfy
\begin{equation}
\AMO(X)\quad\Longleftrightarrow\quad\exists Y\ E(X,Y).
\end{equation}
In particular, an assignment with no true input must remain possible.
All auxiliaries are fresh for each AMO group, including recursive calls.
Groups of size zero or one need no clauses. The formulas below describe
the implementation used to produce the measurements.

\subsubsection{Pairwise}
The direct encoding excludes every pair of true inputs:
\begin{equation}
E_{\mathrm{pair}}(X)=\bigwedge_{1\le i<j\le k}
(\neg x_i\vee\neg x_j).
\end{equation}
Any two true inputs violate their pair clause; conversely, zero or one
true input satisfies all pair clauses. There are $\binom{k}{2}$ clauses
and no auxiliary variables. For $k=3$ these are
$(\neg x_1\vee\neg x_2)$, $(\neg x_1\vee\neg x_3)$, and
$(\neg x_2\vee\neg x_3)$. Setting $x_2$ true immediately forces
$x_1$ and $x_3$ false.

\subsubsection{Binary}
The implementation calls PySAT's bitwise AMO encoder~\cite{pysatcard}.
Let $b=\lceil\log_2 k\rceil$ and introduce $y_0,\ldots,y_{b-1}$.
Assign input $x_i$ the $b$-bit representation of $i-1$, ordered from
most to least significant bit. Write $\ell_{it}=y_t$ when bit position
$t$ in that order is one and $\ell_{it}=\neg y_t$
otherwise. The clauses are
\begin{equation}
E_{\mathrm{bin}}(X,Y)=\bigwedge_{i=1}^{k}\bigwedge_{t=0}^{b-1}
(\neg x_i\vee\ell_{it}).
\end{equation}
Two true inputs require different values of some bit, producing a
contradiction. With one true input, choose its pattern; with none, the
bits are unrestricted. For $k\ge3$, the encoding uses $b$ auxiliaries and
$kb$ clauses. PySAT simplifies $k=2$ to one pairwise clause with no bits.
For four inputs with patterns $00,01,10,11$, selecting $x_3$ fixes pattern
$10$. Each other input disagrees with a fixed bit and is propagated false.

\subsubsection{Commander}
The grouping construction follows Klieber and Kwon~\cite{commander}.
For $k>4$, partition $X$ into consecutive groups $G_1,\ldots,G_m$ of
at most three inputs, where $m=\lceil k/3\rceil$. Introduce commander
$a_j$ for each group and use
\begin{align}
E_{\mathrm{cmd}}(X,A)={}&
\bigwedge_j E_{\mathrm{pair}}(G_j)\ \wedge\
\bigwedge_j\bigwedge_{x\in G_j}(\neg x\vee a_j)\nonumber\\
&{}\wedge\bigwedge_j\left(\neg a_j\vee\bigvee_{x\in G_j}x\right)
\ \wedge\ E_{\mathrm{cmd}}(A).
\end{align}
Recursion ends at $k\le4$ with pairwise clauses. The within-group clauses
exclude two selections in one group; forward implications and commander
AMO exclude selections in different groups. Assigning each commander the
disjunction of its members extends every valid input assignment, including
the all-false assignment. Reverse implications are included in this
implementation, although removing them preserves the projected AMO
relation. This is a stated implementation choice, not a claim that every
commander encoder uses the same clauses or cutoff.

\subsubsection{Recursive product}
This implementation uses the recursive two-dimensional product
construction~\cite{product}. For $k>4$, let $p=\lceil\sqrt{k}\rceil$
and $q=\lceil k/p\rceil$. Introduce $U=(u_1,\ldots,u_p)$ and
$V=(v_1,\ldots,v_q)$. Input $x_i$ occupies
$r(i)=\lfloor(i-1)/q\rfloor+1$, $c(i)=((i-1)\bmod q)+1$. Add
\begin{equation}
E_{\mathrm{prod}}=
\bigwedge_{i=1}^k
\left[(\neg x_i\vee u_{r(i)})\wedge(\neg x_i\vee v_{c(i)})\right]
\wedge E_{\mathrm{prod}}(U)\wedge E_{\mathrm{prod}}(V).
\end{equation}
The recursive AMO calls use pairwise clauses for at most four inputs.
Distinct selected inputs differ in a row or column coordinate and
therefore violate one coordinate AMO. A single selected input can be
extended by choosing its two coordinates; for no selected inputs, choose
all coordinates false. Unused grid cells need no extra exclusion clauses
for this projected AMO relation, and this implementation adds none.
For $k=6$, a $3\times2$ grid maps $x_1,x_2$ to the first grid row,
$x_3,x_4$ to the second, and $x_5,x_6$ to the third.

\subsubsection{Sequential counter}
PySAT supplies the AMO specialization of Sinz's sequential
counter~\cite{sinz}. For $k\ge3$, introduce $s_1,\ldots,s_{k-1}$ and add
\begin{align}
&(\neg x_1\vee s_1),\qquad(\neg x_k\vee\neg s_{k-1}),\nonumber\\
&\bigwedge_{i=2}^{k-1}
\left[(\neg x_i\vee s_i)\wedge(\neg s_{i-1}\vee s_i)
\wedge(\neg x_i\vee\neg s_{i-1})\right].
\end{align}
A true input forces a true prefix indicator that propagates forward,
preventing any later input from being true. A satisfying extension sets
$s_i=\bigvee_{j\le i}x_j$; the clauses do not require this equality for
every auxiliary assignment. The size is $k-1$ auxiliaries and $3k-4$
clauses. PySAT reduces $k=2$ to one pairwise clause without auxiliaries.

\subsubsection{Propagation on the original inputs}
Each construction above has the following useful property for an isolated
AMO: selecting one input makes unit propagation falsify every other input.
Pairwise uses the pair clauses directly. Binary fixes a unique bit pattern.
Commander propagates to the selected group's commander, excludes other
commanders recursively, and then excludes their members. Product selects
one row and column coordinate and excludes every other input through at
least one coordinate. Sequential propagates true prefix indicators forward
and false prefix indicators backward.

For example, with four sequential inputs, setting $x_3$ true forces
$s_3$ true and $s_2$ false. The latter forces $s_1$ and $x_2$ false, then
$x_1$ false; $s_3$ true forces $x_4$ false. This is an input-propagation
property of one AMO, not complete propagation of the entire N-queens CNF.
It also does not imply equal propagation cost or equal search time across
encodings. The proof descriptions follow the displayed clauses; the
additional automated checks are described below.
